\title{Parameter-Efficient Orthogonal Finetuning via Factorization}

\begin{abstract}
The rise of large foundation models has transformed numerous applications, yet training these models from the ground up demands enormous computational resources. Consequently, developing approaches for adapting pre-trained models to specific tasks in a resource-efficient manner has become crucial. This work investigates a systematic adaptation technique, Orthogonal Finetuning (OFT), which facilitates transfer to downstream tasks. While OFT offers strong generalizability, its reliance on high-dimensional orthogonal matrices results in a substantial number of trainable parameters. To mitigate this, we re-examine OFT through the lens of information transfer, identifying several core requirements that lead to improved parameter efficiency. Motivated by the efficiency gains of the Cooley-Tukey fast Fourier transform algorithm, we introduce a novel orthogonal parameterization grounded in butterfly matrix structures. Integrating this parameterization into OFT, we propose Orthogonal Butterfly (BOFT), a new parameter-efficient finetuning strategy. BOFT generalizes OFT by including it as a specific instance, thereby broadening the orthogonal finetuning paradigm. Extensive experimental evaluations are conducted on adapting large vision transformers, large language models, and text-to-image diffusion models across a range of visual and linguistic downstream tasks.
\end{abstract}

\section{Introduction}

Emerging models such as ChatGPT~\cite{brown2020language,chen2021evaluating} and Stable Diffusion~\cite{rombach2022high} exemplify the extraordinary generalization capabilities exhibited by large foundation models. Accompanying these advances in performance is a sharp escalation in model size—

(\emph{e.g.}, GPT-3 encompasses approximately 175B parameters). This trend renders the process of training new foundation models from inception increasingly infeasible for most research groups. As a result, sustained advancement in this area hinges on effective techniques for adapting pre-existing models rather than building new ones from scratch. Specifically, there is a pressing need to develop efficient methods for transferring existing foundation models to address new downstream applications. Three main strategies are typically employed for this adaptation: \emph{model finetuning}~\cite{hulora2022,zaken2022bitfit,guo2020parameter,raffel2020exploring,qiu2023controlling,zhang2022adaptive,chavan2023one}, which retains the original architecture while adjusting only select parameters; \emph{adapter tuning}~\cite{rebuffi2017learning,houlsby2019parameter,pfeiffer2020adapterfusion,lin2020exploring,he2021towards}, involving the integration of additional trainable components into the model; and \emph{prompt tuning}~\cite{li2021prefix,lester2021power}, which modifies the input by appending learnable prefix tokens. Among these alternatives, model finetuning stands out for its simplicity and robustness, and critically, does not introduce additional latency during inference.

The main concept underlying model finetuning is to maintain a close resemblance between the pretrained and finetuned models according to specific metrics, thereby retaining the knowledge acquired during pretraining. Commonly, finetuning techniques utilize a low learning rate, which heuristically constrains the difference between the pretrained and updated models. Recognizing the inherent structure within weight matrices, a more rigorous methodology focuses on safeguarding the relationships within these matrices, specifically preserving the pairwise angular relationships among neurons. This observation motivates the development of Orthogonal Finetuning~(OFT)~\cite{qiu2023controlling}, an innovative approach in which neurons within a layer undergo transformation by a shared orthogonal matrix, ensuring that pairwise angles among neurons remain intact throughout finetuning. While OFT has exhibited strong generalization and convergence properties, particularly for finetuning text-to-image diffusion models~\cite{qiu2023controlling}, it also introduces a substantial parameter count due to the typically large size of orthogonal matrices. To improve efficiency, OFT incorporates a block-diagonal design that decreases the total number of trainable parameters. Nonetheless, this efficiency comes with trade-offs: the imposed block-diagonal pattern restricts the orthogonal transformation within separate blocks, potentially limiting the model’s expressiveness and preventing it from capturing certain traditional linear transformations. Moreover, despite reducing parameter counts, the optimal way to choose a sparsity pattern remains an unresolved issue, and such fixed structures might introduce unfavorable inductive biases.

The central challenge in this context is constructing a dense orthogonal matrix without incurring a prohibitive parameter cost. Ordinarily, forming a dense orthogonal matrix in $d$ dimensions necessitates $\mathcal{O}(d^2)$ parameters, which is often impractical. To overcome this limitation, our methodology diverges from conventional approaches by constructing the dense orthogonal matrix as a product of several factorized sparse matrices. The core idea hinges on the realization that one can reduce the quantity of parameters by accepting increased computation time—essentially, trading space for computational resources. By modeling the orthogonal matrix as the successive product of sparse matrices, each training step requires repeated multiplication, introducing additional computation. 

This matrix factorization process is conceptualized through the lens of information transmission: we interpret the synthesis of a dense orthogonal matrix via sparse matrix products as akin to transmitting information across a graph with a grid structure. Viewing the problem in this way not only elucidates the relationship between sparsity and expressiveness, but also provides the basis for devising various matrix factorization strategies. These strategies are specifically designed to confine the number of learnable parameters while ensuring the ability to form complex, dense matrices.

To realize parameter efficiency within this framework, we draw on the butterfly pattern found in the Cooley-Tukey fast Fourier transform algorithm~\cite{cooley1965algorithm}, which is known for efficiently enabling information exchange between nodes~\cite{klasing1994broadcasting}. This structure naturally leads to what is known as butterfly factorization—a method of sparse matrix factorization guided by the connectivity pattern of the butterfly graph. With butterfly factorization, a $d$-dimensional dense matrix can be expressed as a product of $\mathcal{O}(\log d)$ sparse matrices, where each contains only $\mathcal{O}(d)$ non-zero elements. By maintaining the orthogonality of each sparse factor, the overall orthogonal matrix can be parameterized with merely $\mathcal{O}(d\log d)$ parameters—a dramatic improvement over the traditional dense form. 

Capitalizing on this parameterization, we introduce a new parameter-efficient finetuning approach termed Orthogonal Butterfly~(BOFT). This approach not only generalizes OFT as a particular instance, but also establishes a broad framework for orthogonal finetuning. A fundamental similarity unites both the block-diagonal and butterfly structures: they induce sparsity within orthogonal matrices, thereby reducing the count of adjustable parameters. It is illuminating to juxtapose our method with the low-rank schemes in LoRA~\cite{hulora2022}; both low-rank and sparse matrices are prominent subclasses within the category of structured matrices~\cite{candes2011robust} that excel in achieving parameter efficiency.

In contrast to the block-diagonal architecture utilized by OFT to balance the trade-off between model expressiveness and regularization, BOFT adopts a butterfly structure to create a more gradual transition between matrices drawn from the full orthogonal group (ensuring expressivity) and the identity matrix (promoting regularity). This design facilitates the identification of a more compact hypothesis class within the set of orthogonal matrices for use in downstream applications. Owing to the prevalence of butterfly structures in efficient linear operators—such as the discrete Fourier transform and discrete cosine transform—we posit that this approach induces an implicit inductive bias. This, in turn, enhances model generalization and mitigates overfitting risks. Our reasoning stems from the characteristic that the butterfly structure is capable of expressing various well-known linear transforms, a capability not shared by OFT’s block-diagonal formulation. We outline our main contributions as follows:

\begin{itemize}[leftmargin=*,nosep]
    \item We address the challenge of achieving parameter efficiency in orthogonal finetuning by proposing a new information transmission framework. This approach reframes the problem of constructing dense, parameter-efficient orthogonal matrices as an information flow problem on a grid-based graph.
    \item Drawing motivation from butterfly patterns in the Cooley-Tukey algorithm, we introduce Orthogonal Butterfly, a new orthogonal finetuning technique that emphasizes parameter efficiency within the information transmission paradigm.
    \item We present several theoretical justifications for BOFT’s ability to drastically cut down the number of trainable parameters while maintaining high expressivity and training stability. The factorization property of the butterfly structure in BOFT also enables novel weight interpolation capabilities (refer to Figure~\ref{fig:boft_interpolation}).
    \item For the first time, we extend the application of orthogonal finetuning beyond controllable text-to-image synthesis~\cite{qiu2023controlling}, demonstrating BOFT’s strong suitability as a versatile model finetuning approach. We apply BOFT across a diverse set of adaptation scenarios, including both computer vision and natural language processing tasks. Our results show that BOFT outperforms current leading methods by a substantial margin, thus supporting its advantage in both parameter efficiency and generalization.
\end{itemize}

\section{Related Work}

\textbf{Parameter-efficient finetuning (PEFT)}. With the rapid scaling of foundation models (\emph{e.g.}, \cite{brown2020language,radford2021learning,rombach2022high,kirillov2023segment}), the challenge of adapting these large models to downstream applications using as few parameters as possible has drawn substantial research attention~\cite{treviso2023efficient,ding2023parameter,lialin2023scaling}. A wide range of PEFT approaches have been developed~\cite{houlsby2019parameter,aghajanyan2020intrinsic,hulora2022,edalati2022krona,wang2022adamix,gheini2021cross,zaken2022bitfit,guo2020parameter,sung2021training,ansell2022composable,lester2021power,li2021prefix,vu2022spot,he2021towards,mao2021unipelt,karimi2021compacter,liu2022few,sung2022lst,chen2022parameter,jia2022visual,chen2022adaptformer,zhang2022neural,jie2023fact,lian2022scaling,luo2023towards}, among which those utilizing reparameterization~\cite{aghajanyan2020intrinsic,hulora2022,edalati2022krona} are most pertinent to our investigation. In particular, LoRA~\cite{hulora2022} achieves strong results on various natural language processing tasks by augmenting pretrained weights with the product of two learnable low-rank matrices. LoRA prescribes a uniform rank for all inserted matrices, leading to subsequent efforts~\cite{zhang2022adaptive,valipour2022dylora,zhang2023increlora} that dynamically tailor the rank per layer, enabling more effective parameter distribution. The simplicity of this low-rank strategy has driven widespread adoption~\cite{dettmers2023qlora,zi2023delta,chavan2023one}. Inspired by the role of hyperspherical energy in describing generalization~\cite{liu2018learning,liu2021orthogonal}, \cite{qiu2023controlling} introduce orthogonal finetuning (OFT) as an alternative technique for fine-tuning text-to-image diffusion models. OFT involves optimizing an orthogonal transformation within each layer, resulting in superior generalization and training stability compared to LoRA. Nonetheless, OFT typically incurs a larger parameter overhead relative to LoRA. As such, advancing OFT’s parameter efficiency is both valuable and necessary. Furthermore, OFT’s effectiveness outside its original domain (text-to-image diffusion control~\cite{qiu2023controlling}) remains untested. BOFT addresses OFT's inefficiency by leveraging butterfly factorization, which now enables a broader demonstration of orthogonal finetuning’s adaptability across different tasks.

\textbf{Butterfly structures}. The radix-2 Cooley-Tukey scheme decomposes the $N$-point discrete Fourier transform into smaller $\frac{N}{2}$-point transforms recursively, naturally giving rise to the so-called butterfly structure—represented as a sequence of sparse matrix multiplications, often referred to as a butterfly matrix. Butterfly matrices have already found use in parameterizing orthogonal matrices to circumvent pivoting in Gaussian elimination and boost computational efficiency~\cite{parker1995random}; in recurrent neural networks for improving training stability~\cite{jing2017tunable}; and in kernel approximation contexts~\cite{munkhoeva2018quadrature}. Additionally, \cite{dao2019learning,dao2019kaleidoscope} exploit butterfly parameterizations to learn fast linear operators, while \cite{chen2022pixelated,dao2022monarch} apply butterfly matrices for efficient neural network training. These structures have proven their utility for accelerating matrix-vector multiplication~\cite{michielssen1996multilevel,o2010algorithm}, for compact matrix representation~\cite{li2015butterfly}, and for communication network protocols~\cite{klasing1994broadcasting,li2003linear,rajasingh2016transmission}. In this work, diverging from these applications, our objective is to utilize butterfly structures for boosting the parameter efficiency of OFT.

\section{An Information Transmission View on Orthogonal Finetuning}

We begin by outlining some foundational concepts of OFT. In order to adapt a pretrained weight matrix, OFT introduces a new parameterization where the updated weights are expressed as the product of a trainable orthogonal matrix and the original (fixed) pretrained matrix. Unlike LoRA, which performs adaptation by adding a low-rank matrix to the original weights, OFT employs a multiplicative update using an orthogonal transformation. LoRA achieves parameter efficiency via a low-rank parameterization, whereas the standard OFT leverages a block-diagonal structure for its orthogonal matrix~\cite{qiu2023controlling}; reducing the block size within the diagonal blocks increases the parameter efficiency of OFT. Figure~\ref{fig:comp_lora} visually summarizes these differences. 

The key rationale behind utilizing orthogonal transformations during finetuning is to maintain the pairwise angular relationships between neurons~\cite{liu2018learning,liu2021orthogonal,qiu2023controlling}, which helps retain the semantic information acquired during pretraining. Specifically, OFT optimizes an orthogonal matrix $\bm{R}\in\mathbb{R}^{d\times d}$ applied to a pretrained linear mapping $\bm{W}^0\in\mathbb{R}^{d\times n}$, thereby altering the usual forward computation from $\bm{z}=(\bm{W}^0)^\top\bm{x}$ to $\bm{z}=(\bm{R}\bm{W}^0)^\top\bm{x}$, where $\bm{x}\in\mathbb{R}^{d}$ represents the input and $\bm{z}\in\mathbb{R}^{n}$ the output. To initialize with zero perturbation, OFT sets $\bm{R}$ to the identity matrix. To keep $\bm{R}$ orthogonal throughout finetuning, we adopt the Cayley parameterization~\cite{qiu2023controlling,liu2021orthogonal}, namely $\bm{R}=(\bm{I}+\bm{Q})(\bm{I}-\bm{Q})^{-1}$, with $\bm{Q}$ constrained to be skew-symmetric, i.e., $\bm{Q}=-\bm{Q}^\top$.

For further efficiency, the block-diagonal approach parameterizes the orthogonal matrix as $\text{diag}(\bm{R}_1,\bm{R}_2,\cdots,\bm{R}_r)$, where each $\bm{R}_i\in\mathbb{R}^{b\times b}$ is a small orthogonal matrix, and $br=d$. The use of a block-diagonal form enhances parameter efficiency but introduces a structural assumption: the neuron's dimensionality (the columns of $\bm{W}^0$) is segmented into $r$ groups, each transformed by its own orthogonal block. Although this blockwise partitioning yields practical gains, it is arbitrary to group a neuron's dimensions merely by their index, thus highlighting the limitations of the block-diagonal strategy and motivating the exploration of dense orthogonal matrices. This brings us to a pivotal question: \emph{Is it possible to design a dense orthogonal matrix while still maintaining parameter efficiency?}

In response to this issue, we suggest constructing a dense orthogonal matrix by multiplying several sparse orthogonal matrices together. To analyze the structure of sparsity in orthogonal matrix decomposition from a coherent and accessible standpoint, we recast the generation of a dense orthogonal matrix in OFT as an instance of information transmission. Concretely, forming a dense matrix $\bm{R}\in\mathbb{R}^{d\times d}$ as the product of $m$ square matrices $\thickmuskip=2mu \medmuskip=2mu \bm{R}=\bm{B}_m\bm{B}_{m-1}\cdots\bm{B}_1$ is analogous to conveying information over a grid consisting of $d\times (m+1)$ nodes, as depicted in Figure~\ref{fig:info_trans}. This perspective is motivated by the observation that a dense $d$-dimensional square matrix embodies full connectivity between two sets of $d$ nodes. For the matrix $\bm{R}$, a zero entry $R_{ij}$ signifies that no information can be transmitted from the $j$-th source node to the $i$-th destination node, whereas a nonzero $R_{ij}$ allows such transmission. Therefore, decomposing $\bm{R}$ into multiple matrices $\bm{B}_m\bm{B}_{m-1}\cdots\bm{B}_1$ can be viewed as a sequence of information exchanges along the sequence of graphs determined by each $\bm{B}_i$. Information is first relayed according to $\bm{B}_1$ and ultimately transferred via $\bm{B}_n$. For example, Figure~\ref{fig:info_trans} illustrates the case of $\bm{R}=\bm{B}_5\bm{B}_4\bm{B}_3\bm{B}_2\bm{B}_1$, with their respective sparsity patterns and resulting graph structure displayed. This graph reflects the process of unrolling the matrix product. Here, each $\bm{B}_i$ acts as a connectivity matrix linking nodes at the $i$-th level to those at the $(i+1)$-th level. In more detail, the $(j_1,j_2)$ entry in $\bm{B}_i$ indicates whether there exists a directed edge from the $j_2$-th node at level $i$ to the $j_1$-th node at level $(i+1)$ (with a value of zero corresponding to no connection). For the product $\bm{B}_5\bm{B}_4\bm{B}_3\bm{B}_2\bm{B}_1$ to yield a dense matrix, it is required that each node in the first layer can send information to all nodes in the sixth layer. Examining only the product $\bm{R}=\bm{B}_4\bm{B}_3\bm{B}_2\bm{B}_1$—connecting the source at level one to the recipients at level five—reveals that node $1$ cannot communicate with node $3$, meaning the $(3,1)$ entry in the resulting sparsity pattern is zero. Analogously, a similar relationship between the induced graphs and sparsity patterns can be observed for $\bm{R}=\bm{B}_3\bm{B}_2\bm{B}_1$ and $\bm{R}=\bm{B}_2\bm{B}_1$. More generally, to achieve a dense $\bm{R}\in\mathbb{R}^{d\times d}$, the collection of factor matrices $\bm{B}_m,\cdots,\bm{B}_1$ must represent a set of directed edges over a $d\times (m+1)$ node grid, where each edge only connects nodes at adjacent layers (i.e., adjacent columns), guaranteeing that information originating from any first-level node reaches every node at the $(m+1)$-th level.

The perspective of information transmission offers valuable insight into the sparsity characteristics of the factorization matrices used in OFT. In Figure~\ref{fig:blk_diag}, the block-diagonal nature of $\bm{R}$ in the original OFT is illustrated. While this design decreases the overall number of parameters that require training, the block-diagonal structure limits the possibility of forming a dense matrix $\bm{R}$. The objective, therefore, is to generate a dense orthogonal matrix by combining $m$ sparse orthogonal matrices, all while minimizing the number of effective trainable parameters. From the information transmission perspective, two primary criteria are desirable: \emph{(i)} \textbf{dense connectivity}—ensuring each node in the input layer can reach any node in the output layer through at least one path; and \emph{(ii)} \textbf{minimal free edges}—maintaining the smallest possible edge count while upholding orthogonality. It is important to recognize that orthogonality imposes a stringent requirement on the connections across adjacent layers. For instance, maintaining full-rank property in each $\bm{B}_i$ (essential for orthogonality) necessitates at least $d$ edges to create a bijective correspondence between the $i$-th and $(i=1)$-th layers, ensuring no fewer than $d$ connections between layers (as exemplified by the 4-edge scenario in Figure~\ref{fig:info_trans}). These $d$ connections, required solely to achieve orthogonality, are not counted among trainable edges since their values are fixed and non-trainable (\emph{e.g.}, for a $d\times d$ orthogonal matrix containing $d$ nonzero elements, each entry must be either $+1$ or $-1$). Because the orthogonality condition introduces additional dependencies across the edges, fewer trainable parameters are necessary compared to unrestricted full matrices. For example, in the case of a $2\times 2$ orthogonal matrix, a zero entry at position $(1,1)$ mandates a corresponding zero at position $(2,2)$ (\emph{i.e.}, one absent edge enforces the absence of another). Hence, the structure induced by orthogonality can lead to either the addition or removal of edges, depending on the specific connection pattern. Focusing on the arrangement of nonzero elements within the composed orthogonal matrix, the information transmission viewpoint is especially advantageous for OFT, since only the placement of trainable nonzero elements in $\bm{R}$ is relevant—the actual values they assume are inconsequential.

A straightforward fully-connected mapping between two layers requires $\mathcal{O}(d^2)$ connections—equivalent to a dense orthogonal matrix—resulting in $d^2-d$ actual edges (for instance, with $d=4$, this amounts to 12 edges). In Figure~\ref{fig:info_trans}, a concrete example is presented showing a viable matrix factorization, which utilizes a total of $10$ edges, fewer than needed for a conventional dense orthogonal matrix. Utilizing this construction allows us to assess the parameter efficiency of OFT graphically, and we can systematically devise valid factorizations under this paradigm. We further take cues from the unique network topology found in the Cooley-Tukey algorithm, namely the butterfly graph~\cite{cooley1965algorithm}, which is known to efficiently establish dense connections between $d$ input and $d$ output nodes using only $\mathcal{O}(d\log d)$ edges. To illustrate, while the configuration in Figure~\ref{fig:info_trans} achieves dense connections with $10$ edges, the butterfly network achieves the same with just $8$. In what follows, we discuss how leveraging the butterfly topology enhances parameter efficiency.

\section{Orthogonal Parameterization by Butterfly Factorization}

The butterfly topology originally emerged as a key component in the Cooley-Tukey algorithm for the fast Fourier transform. Within the Fourier transform, a minor alteration in frequency can propagate to a widespread modification in the spatial domain. This property mirrors our concern with information flow, as we also seek for all nodes in the initial layer to access all nodes in the final layer. The butterfly network has also gained traction in the design of computer architectures for optimizing data communication~\cite{solihin2015fundamentals,li2011linear}. Assuming $k\geq 2$ is a power of $2$, we define the so-called \emph{butterfly factor} $\bm{B}^F(k)$ as:

\begin{equation}\label{eq:but_factor}
\begin{aligned}
    \bm{B}^F(k)=\begin{bmatrix}
    \text{diag}(\bm{d}_1) & \text{diag}(\bm{d}_2) \\
    \text{diag}(\bm{d}_3) & \text{diag}(\bm{d}_4)
    \end{bmatrix}\in\mathbb{R}^{k\times k},
\end{aligned}
\end{equation}

where each vector $\bm{d}_i\in\mathbb{R}^{\frac{k}{2}}$ for $i=1,2,3,4$. For $d=2^N$, we then recursively construct the $d$-dimensional \emph{butterfly matrix} $\bm{B}(d)\in\mathbb{R}^{d\times d}$ as:

\begin{equation}
\begin{aligned}
    \!\!\bm{B}(d)=\tilde{\bm{B}}(d,d)\!\cdot\!\begin{bmatrix}
    \bm{B}_1(\frac{d}{2})\!\!\!\!\! & \bm{0} \\
    \bm{0}\!\!\!\!\! &\bm{B}_2(\frac{d}{2})
    \end{bmatrix}= \tilde{\bm{B}}(d,d)\tilde{\bm{B}}(d,\frac{d}{2})\cdots\tilde{\bm{B}}(d,2),
\end{aligned}
\end{equation}

Here, $\bm{B}_1(\frac{d}{2})$ and $\bm{B}_2(\frac{d}{2})$ denote two butterfly matrices, each of dimension $\frac{d}{2}$. Next, we introduce the \emph{butterfly component}, given by $\thickmuskip=2mu \medmuskip=2mu \tilde{\bm{B}}(d,k)=\text{diag}(\bm{B}^F_1(k),\cdots,\bm{B}^F_{d/k}(k))$, which is constructed as a block-diagonal matrix of dimension $d\times d$, where each diagonal block is of size $k$ and every $\bm{B}^F_i(k)$ represents a butterfly factor as specified in Equation~\ref{eq:but_factor}. Utilizing this construction, we aim to parameterize an orthogonal matrix through the butterfly matrix. This is accomplished by requiring that each multiplicative block $\tilde{\bm{B}}(d,k)$ within the butterfly structure $\bm{B}(d)$ is itself orthogonal. Focusing first on the special case where the block size is $2$, i.e., the matrix $\tilde{\bm{B}}(d,2)$, orthogonality can be simply achieved by employing the Cayley transform or implementing $2$-dimensional rotations for each block, as outlined in \cite{liu2021orthogonal,qiu2023controlling}. The nonzero structure of any butterfly component reflects a permutation of the nonzero pattern present in $\tilde{\bm{B}}(d,2)$, which allows us to efficiently ensure orthogonality across all butterfly components via analogous parameterization. This approach results in a scalable method for constructing orthogonal matrices by composing multiple $2\times 2$ orthogonal blocks. Following the generalization in \cite{chen2022pixelated}, we define a block-wise butterfly component $\tilde{\bm{B}}^b(d,k)$, where every element of $\bm{d}_i$ for all $i$ is replaced with a $b\times b$ matrix. To ensure that the block butterfly component $\tilde{\bm{B}}^b(d,2)$ remains orthogonal, we parameterize each $2b\times 2b$ block to be orthogonal. For other butterfly components $\tilde{\bm{B}}^b(d,k)$ with $k>2$, the nonzero block patterns are derived through permutation from the pattern in $\tilde{\bm{B}}^b(d,2)$, so orthogonality can be enforced in the same way. Collectively, these mechanisms facilitate an efficient forward pass within BOFT.

\begin{equation}
    \begin{aligned}\nonumber
    \bm{z}=\big(\bm{R}(m,b)\cdot\bm{W}^0\big)^\top\bm{x},~~~~\text{s.t.}~\bigg{\{}\bm{R}(m,b)=\prod_{i=1}^m \tilde{\bm{B}}^b_{(d,i)}~~\&~~\underbrace{\big(\tilde{\bm{B}}^b_{(d,j)}\big)^\top\tilde{\bm{B}}^b_{(d,j)}=\tilde{\bm{B}}^b_{(d,j)}\big(\tilde{\bm{B}}^b_{(d,j)}\big)^\top\!=\bm{I}_d}_{\forall j\in [1, m]}\bigg{\}},
    \end{aligned}
\end{equation}

In this setting, the notation $\tilde{\bm{B}}^b(d,2^{m - i+1})$ is abbreviated as $\tilde{\bm{B}}^b_{(d,i)}$ for clarity, and $\bm{I}_d$ indicates the $d \times d$ identity matrix. The orthogonal matrix $\bm{R}(m,b)\in\mathbb{R}^{d\times d}$ is structured as a sequential product of $m$ orthogonal butterfly factors. To simplify references, we write BOFT with $\bm{R}(m,\frac{b}{2})$ as BOFT($m$,$b$) where $b\geq 2$. When the parameters satisfy $m=1$, BOFT($1$,$b$) corresponds to the block-diagonal OFT~\cite{qiu2023controlling} configuration with block size $b$. For the particular choice of BOFT($1$,$d$), the approach collapses to the canonical OFT~\cite{qiu2023controlling} comprising a fully unconstrained orthogonal matrix. Moreover, selecting BOFT($\log \frac{2d}{b}$,$b$) amounts to using the block butterfly matrix $\bm{B}^{\frac{b}{2}}(d)$ as $\bm{R}$, which produces a dense orthogonal matrix structure. Broadly, the total count of effective tunable parameters in BOFT($m$,$b$) scales as $\frac{1}{2}(b-1)dm$ for fine-tuning a linear transformation of dimensions $d\times n$. In the case where a butterfly matrix is applied, specifically with $m=\log d, b=2$, the parameter count becomes $\mathcal{O}(d\log d)$. This is in contrast with the original OFT with an unconstrained orthogonal matrix, which has $\mathcal{O}(d^2)$ parameters, or the block-diagonal variant with $r$ blocks requiring $\mathcal{O}(bd)$. Consequently, the original OFT needs $b=d$ to yield a fully dense orthogonal matrix, whereas BOFT provides such flexibility for any block size $b$.

\textbf{Identity initialization for BOFT}. In standard fine-tuning pipelines, initializing from the original pretrained weights is common practice to maintain proximity between the new and pretrained models. As an instance, LoRA chooses to set the low-rank component's weights to zero at initialization. Analogously, BOFT starts by initializing all butterfly factors as identity matrices (equivalently, the skew-symmetric matrix in the Cayley transform is set to zero).

\textbf{Multiplicative Dropout for BOFT}. LoRA~\cite{hulora2022} introduces a Dropout layer on the low-rank parameter update to minimize overfitting risk. While classical Dropout~\cite{srivastava2014dropout} is compatible with LoRA's additive model, it is not naturally suited to the multiplicative structure employed by BOFT. To overcome this limitation, we devise a multiplicative Dropout mechanism tailored for BOFT. Notably, since $\bm{R}(m,b)$ is a product of $m$ orthogonal butterfly factors that can each be reconfigured into $2b\times 2b$ block-diagonal orthogonal forms, our Dropout procedure randomly selects $p_1$ percent of butterfly factors and $p_2$ percent of the diagonal blocks within each, replacing these selected blocks with identity matrices during training.

\section{Intriguing Insights and Discussions}

\textbf{Expressivity of BOFT}. The butterfly architecture, combined with permutations, is known to exactly recover a range of well-established fast linear transforms~\cite{dao2019learning,dao2019kaleidoscope} (\emph{e.g.}, the fast Fourier transform and the Hadamard transform). However, the degree to which our orthogonal butterfly matrices can approximate an arbitrary orthogonal matrix has not been fully characterized. To investigate this, we perform a numerical experiment aimed at approximating a randomly generated, dense orthogonal matrix~\cite{anderson1987generation} of dimension $512\times 512$. The outcomes, visualized in Figure~\ref{fig:para_eff}, reflect averages over 10 independently drawn random seeds. On the plot, the y-axis measures the approximation error, while the x-axis represents the count of effective trainable parameters. For each block size, BOFT’s results are depicted in the same color; the leftmost datapoint of each color shows the error for BOFT($1$,$b$) (which reduces to the canonical block-diagonal OFT with block size $b$). Across the board, BOFT achieves superior parameter efficiency relative to OFT. For instance, BOFT($9$,$2$) attains a better approximation than BOFT($1$,$16$), despite requiring fewer parameters. Generally, adopting a smaller $b$ and larger $m$ for BOFT leads to enhanced parameter efficiency; e.g., BOFT($6$,$4$) consumes significantly fewer parameters while achieving an error comparable to that of BOFT($2$,$16$). Collectively, these observations highlight that the butterfly structure parametrizes a more restrictive yet expressive subset of the orthogonal group than the block-diagonal alternative, thereby making BOFT provably more expressive than OFT when block sizes are matched.

\begin{theorem}[Expressivity of BOFT]\label{thm:exp_boft}
    BOFT possesses greater expressivity than OFT under equal block sizes. Any orthogonal matrix of dimension $d$ can, in theory, be constructed via products of butterfly matrices in the form $\bm{B}_{d-1,1}(d)\bm{B}^\top_{d-1,2}(d)\cdots\bm{B}_{1,1}(d)\bm{B}^\top_{1,2}(d)$, with each $\bm{B}_{i,j}(d)$ denoting a butterfly matrix for any $i$ and $j$.
\end{theorem}

The statement in Theorem~\ref{thm:exp_boft} motivates a straightforward extension for BOFT: specifically, the target orthogonal matrix can be represented as $\thickmuskip=2mu \medmuskip=2mu \bm{R}^G(m_1,b_1,m_2,b_2,l)=\bm{R}_{l,1}(m_1,b_1)\bm{R}_{l,2}^T(m_2,b_2)\cdots\bm{R}_{1,1}(m_1,b_1)\bm{R}_{1,2}^T(m_2,b_2)$, where $\bm{R}_{i,j}^T(m,b)$ refers to the orthogonal factors within BOFT. By selecting $m_1=m_2=\log d$, $b_1=b_2=2$, and $l=d-1$, the resulting composition $\bm{R}^G(m_1,b_1,m_2,b_2,l)$ spans the entire orthogonal group, a construction also termed the kaleidoscope hierarchy~\cite{dao2019kaleidoscope}. It should be noted, however, that augmenting expressiveness is not guaranteed to improve performance during finetuning. For instance, unconstrained full finetuning (with its maximal representational power) often underperforms in practice. Ultimately, balancing the competing demands of expressivity and regularization determines the success of model finetuning. By embedding structural priors, BOFT effectively enlarges the parameter space available for finetuning, allowing for a more optimal balance between expressivity and regularity.

\textbf{Spectral properties}. Compared to LoRA, orthogonal finetuning more effectively maintains the spectral characteristics of the pretrained weight matrix $\bm{W}^0$, as it exactly conserves the spectral norm. This can be demonstrated by performing singular value decomposition: $\bm{W}^0=\bm{U}\bm{\Sigma}\bm{V}^\top$, where $\bm{U}$ and $\bm{V}$ denote orthogonal matrices and $\bm{\Sigma}$ is a diagonal matrix of singular values. In both OFT and BOFT, an orthogonal transformation $\bm{R}$ is applied to the left, resulting in finetuned parameters $\bm{R}\bm{U}\bm{\Sigma}\bm{V}^\top$. This operation leaves the dominant singular value, and thus the spectral norm of $\bm{W}^0$, unaffected. Preserving the spectral norm in this way has been shown to substantially improve both the stability of training and the model's generalization ability~\cite{miyato2018spectral,yoshida2017spectral}. Further mathematical discussion can be found in Appendix~\ref{app:sn_property}.

\textbf{Orthogonal finetuning as learning bilinear similarity}. The BOFT approach can be reformulated as optimizing a bilinear similarity function of the form $\bm{w}^0_i\bm{R}\bm{x}$, where $\bm{w}^0_i$ denotes the $i$-th column (neuron) in $\bm{W}^0$. This formulation interprets BOFT as the task of learning a bilinear similarity transformation matrix $\bm{R}$ under a strict orthogonality constraint, directly relating to the fields of distance metric learning~\cite{xing2002distance} and bilinear forms~\cite{roman2005advanced}.

\textbf{Inductive bias and generalization in BOFT}. In BOFT, the transformation $\bm{R}(m,b)$ typically belongs to a specially structured subset of the orthogonal group, thereby restricting the expressivity of the hypothesis space and introducing an inductive bias. We propose that this inductive bias—imposed through butterfly factorization—is advantageous for generalization, since it encodes structural patterns common to well-known linear operators~\cite{dao2019kaleidoscope}, such as the discrete Fourier, discrete sine/cosine, and Hadamard transforms. Additionally, the sparsity induced by matrix factorization in BOFT can introduce further forms of implicit bias~\cite{gunasekar2017implicit,li2018algorithmic,liu2019neural}.

\textbf{Comparison to butterfly-based sparse training}. Several studies~\cite{dao2019learning,dao2019kaleidoscope,chen2022pixelated,dao2022monarch} have explored the use of butterfly parameterization in sparse training. These works predominantly reparameterize weight matrices using the butterfly parameterization and conduct end-to-end training of neural networks from initialization. \cite{dao2022monarch} examines an approach where pretrained weights are first projected onto modified butterfly matrices, after which the projected segments are further optimized for specific downstream tasks. In contrast, BOFT introduces a substantially different finetuning methodology by modifying the weights with weight matrices that are shared across layers.

\input{rebuttal-results}

\section{Concluding Remarks and Limitations}

In this work, we introduce Orthogonal Butterfly, a parameter-efficient, general-purpose finetuning technique for foundation models utilizing the butterfly structure. Our primary insight for enhancing parameter-efficiency is to express a dense orthogonal matrix as the product of several sparse orthogonal matrices. To systematically identify valid factorizations, we put forward a graphical information transmission framework. Within this framework, we observe that the butterfly structure reliably meets the requirements for sparse orthogonal matrix factorization. The empirical studies conducted confirm the practical effectiveness of BOFT across diverse tasks, including finetuning large language models, large-scale vision models, and text-to-image generative models. Our results further establish BOFT’s advantages as a broadly applicable finetuning approach.

However, BOFT is not without its drawbacks. As the resultant orthogonal matrix is built by multiplying multiple orthogonal matrices, the associated training runtime is somewhat higher than that of OFT. Enhancing the training efficiency of BOFT is an open area for future work. On the upside, once finetuning is complete, the orthogonal matrices learned by BOFT can be merged directly into the pretrained model’s parameters, incurring no extra inference time. Additionally, it remains an open question whether the butterfly network offers the optimal approach for information transmission. Our information transmission framework also enables us to look to other fields—such as computer networking, where the study of efficient topologies for information transfer is well-established—for further inspiration. We anticipate that network architectures even more effective than the butterfly structure may be devised for constructing dense orthogonal matrices.

\end{document}